\section{Jacobi rotation method}

\subsection{Task}

Implement the rotation method taking the matrix and the precision as input. Find the eigenvalues
and the eigenvectors of the symmetric matrix and analyze the dependence of the error on the number
of iterations.

\[
A = \begin{pmatrix}
3 & 2 & 6 \\
2 & -3 & -7 \\
6 & -7 & 3
\end{pmatrix}
\]

\subsection{Method}

The method builds a sequence of similar matrices $A^{(k+1)} = U^{(k)T} A^{(k)} U^{(k)}$, where
$U^{(k)}$ is the rotation in the plane $(i, j)$ of the largest off-diagonal element
$|a_{ij}^{(k)}|$:
$u_{ii} = u_{jj} = \cos\varphi_k$, $u_{ij} = -\sin\varphi_k$, $u_{ji} = \sin\varphi_k$.
The angle makes the element $a_{ij}^{(k+1)}$ zero:

\[
\varphi_k = \frac{1}{2} \arctan \frac{2 a_{ij}^{(k)}}{a_{ii}^{(k)} - a_{jj}^{(k)}},
\qquad
\varphi_k = \frac{\pi}{4} \operatorname{sign} a_{ij}^{(k)} \text{ if } a_{ii}^{(k)} = a_{jj}^{(k)}.
\]

A rotation changes only rows and columns $i$ and $j$, so one step costs $O(n)$ plus $O(n^2)$ for
the search of the pivot. The error is measured by the norm of the off-diagonal part

\[
t\bigl(A^{(k)}\bigr) = \Bigl(\sum_{i<j} \bigl(a_{ij}^{(k)}\bigr)^2\Bigr)^{1/2},
\]

and the process stops when $t(A^{(k)}) \le \varepsilon$. The diagonal of $A^{(k)}$ contains the
eigenvalues and the columns of $U = U^{(0)} U^{(1)} \cdots U^{(k-1)}$ are the orthonormal
eigenvectors.

\subsection{Results}

The program reads $n$, the matrix $A$ and the precision $\varepsilon$ and prints the error after
every rotation, the eigenvalues and the eigenvectors.

\needspace{10\baselineskip}
\lstinputlisting[style=console, caption={tests/data/rotation/01.in}]{../tests/data/rotation/01.in}
\lstinputlisting[style=console, caption={tests/data/rotation/01.out}]{../tests/data/rotation/01.out}

The unit tests check $A x_i = \lambda_i x_i$ and $U^T U = E$ on random symmetric matrices up to
$7 \times 7$.

\subsection{Analysis}

A rotation preserves the Frobenius norm and moves the weight of the element $a_{ij}$ to the
diagonal, so $t^2(A^{(k+1)}) = t^2(A^{(k)}) - \bigl(a_{ij}^{(k)}\bigr)^2$ and the error decreases
monotonically (the first step of the table removes $a_{23} = -7$: $9.434^2 - 6.325^2 = 49$).
Since the largest of the $n(n-1)/2$ elements is chosen, $a_{ij}^2 \ge \frac{2}{n(n-1)} t^2$, and
the convergence is at least linear:

\[
t\bigl(A^{(k)}\bigr) \le \Bigl(1 - \frac{2}{n(n-1)}\Bigr)^{k/2} t\bigl(A^{(0)}\bigr).
\]

For $n = 3$ the guaranteed factor is $\sqrt{2/3} \approx 0.82$ per rotation, but the table shows
a faster decrease. The first four rotations reduce the error by 1.5--2.7 times, and then the
convergence becomes quadratic:
$4.9 \cdot 10^{-2} \to 3.3 \cdot 10^{-3} \to 2.1 \cdot 10^{-5} \to 3.4 \cdot 10^{-9}$,
the number of correct digits roughly doubles with every rotation. That is why 8 rotations are
enough for $\varepsilon = 10^{-6}$, and a higher precision costs only one or two extra rotations.

\subsection{Source code}

\lstinputlisting[caption={include/rotation.hpp}]{../include/rotation.hpp}
\lstinputlisting[caption={src/rotation.cpp}]{../src/rotation.cpp}

\pagebreak
